Software verification is the discipline of establishing that a program meets its
specification and is indispensable to improve our understanding of software 
systems and ensure their reliability. 
Theorem proving pursues verification by constructing formal machine-checkable 
proofs of correctness, in a logic that can capture program semantics.
Constructing such proofs turns repeatedly on deciding whether terms are equal,
and many theorem provers rest on a common foundation for equational reasoning.
Reasoning in verification, however, is rarely about equality alone. 
Proofs split into cases, argue under assumptions, and appeal to facts that hold 
only conditionally, but provers today capture this richer semantics around that 
foundation rather than within it. This proposal argues that the foundation
itself should reason beyond equality, letting provers do more of their work at
the level they already share.

We pursue this idea for the e-graph, the equality-reasoning data structure at 
the heart of many provers. Our project, \system, moves that richer semantics
inside the e-graph: disequalities and conditional facts become part of the
structure itself, and provers reason about them natively. These semantic
extensions open room for optimizations, since one enriched structure can reuse
work that separate encodings would repeat, and because the e-graph underlies so
many provers, these gains extend well beyond any single tool. We plan to realize our
idea in a single artifact that integrates with existing provers and, through
them, reaches the broader verification ecosystem.